\subsection*{قسمت ۱: تعداد پارامترهای لایه‌ی آخر}

ابتدا ابعاد فضایی خروجی هر لایه را دنبال می‌کنیم تا به ورودی لایه‌ی تمام‌متصل برسیم. فرمول ابعاد خروجی:
$$O = \left\lfloor \frac{I - K + 2P}{S} \right\rfloor + 1$$

برای \lr{Conv1} با ورودی $224$، $K=7$، $P=3$، $S=2$:
$$O_1 = \left\lfloor \frac{224 - 7 + 6}{2} \right\rfloor + 1 = \left\lfloor \frac{223}{2} \right\rfloor + 1 = 111 + 1 = 112$$
پس خروجی $112 \times 112 \times 64$ است.

برای \lr{Conv2} با $K=3$، $P=1$، $S=1$:
$$O_2 = \frac{112 - 3 + 2}{1} + 1 = 112$$
خروجی $112 \times 112 \times 128$.

برای \lr{MaxPool} با $K=3$، $S=2$، بدون \lr{padding}:
$$O_3 = \left\lfloor \frac{112 - 3}{2} \right\rfloor + 1 = 54 + 1 = 55$$
خروجی $55 \times 55 \times 128$.

برای \lr{Conv3} با $K=3$، $P=1$، $S=2$:
$$O_4 = \left\lfloor \frac{55 - 3 + 2}{2} \right\rfloor + 1 = 27 + 1 = 28$$
خروجی $28 \times 28 \times 256$.

بنابراین بردار ورودی لایه‌ی تمام‌متصل دارای $28 \times 28 \times 256 = 200704$ نورون است و خروجی $10$ نورون. با احتساب \lr{bias}:
$$\#params = 200704 \times 10 + 10 = 2{,}007{,}050$$

\subsection*{قسمت ۲: کاهش پارامترها با کانولوشن عمق‌جدایی‌پذیر}

در \lr{Conv3} استاندارد، ورودی $128$ کانال و خروجی $256$ کانال با فیلتر $3\times 3$ است؛ پس:
$$P_{std} = 3 \times 3 \times 128 \times 256 = 294912$$

در حالت \lr{Depth-wise Separable}، ابتدا یک \lr{depthwise} روی هر کانال ورودی و سپس یک \lr{pointwise} ($1\times1$) برای ترکیب کانال‌ها:
$$P_{dw} = 3 \times 3 \times 128 = 1152, \qquad P_{pw} = 1 \times 1 \times 128 \times 256 = 32768$$
$$P_{ds} = 1152 + 32768 = 33920$$

از این رو میزان کاهش:
$$\frac{P_{ds}}{P_{std}} = \frac{33920}{294912} \approx 0.115$$
یعنی تعداد پارامترها حدود $8.7$ برابر کم می‌شود (به حدود $11.5\%$ مقدار اولیه می‌رسد).

\subsection*{قسمت ۳: عملکرد ضعیف روی اشیای حاشیه‌ای}

\subsubsection*{آ}
وجود لایه‌ی تمام‌متصل در انتهای شبکه باعث می‌شود مکان قرارگیری شیء در تصویر در آموزش کدگذاری شود؛ چون داده‌ی آموزشی عمدتاً اشیا را در مرکز تصویر دارد، شبکه یاد می‌گیرد فعال‌سازی‌های مربوط به مرکز را به کلاس‌ها نگاشت کند و برای موقعیت‌های نامعمول تعمیم نمی‌یابد. این سوگیری مکانی (نه ضعف استخراج ویژگی) عامل اصلی است.

\subsubsection*{ب}
لایه‌های کانولوشن \lr{Translation Equivariant} هستند، یعنی جابه‌جایی ورودی به جابه‌جایی متناظر در نقشه‌ی ویژگی منجر می‌شود؛ اما این خاصیت تنها تا قبل از لایه‌ی \lr{FC} حفظ می‌شود. لایه‌ی تمام‌متصل، نقشه‌ی ویژگی را به‌صورت تخت (\lr{flatten}) و هر موقعیت را به وزن متفاوتی متصل می‌کند؛ پس کل شبکه \lr{Translation Invariant} نیست و به مکان حساس می‌شود. به‌علاوه \lr{padding} مرزها نیز خاصیت تغییرناپذیری مکانی را در حاشیه‌ها می‌شکند و سرنخ مکانی به شبکه می‌دهد.

\subsubsection*{ج}
بدون تغییر معماری، می‌توان با \lr{Data Augmentation} مکانی (جابه‌جایی تصادفی، \lr{random crop}، \lr{translation}) داده‌ها را غنی کرد تا اشیا در موقعیت‌های متنوع دیده شوند و سوگیری مرکزی از بین برود.

\section*{پرسش ۴}

خروجی \lr{Conv3} برابر $28 \times 28 \times 256$ است و \lr{ROI-Pooling} نتیجه را به ابعاد ثابتی نگاشت می‌کند که پس از تخت‌شدن باید با ورودی $1024$ بُعدی لایه‌ی اول سازگار باشد. اگر خروجی \lr{ROI-Pooling} ابعاد $H \times W \times 256$ داشته باشد، آنگاه:
$$H \times W \times 256 = 1024 \implies H \times W = 4$$
پس $H = W = 2$؛ یعنی \lr{ROI-Pooling} باید طول و عرض $2 \times 2$ داشته باشد.

\section*{پرسش ۵}

دلیل اصلی نیاز \lr{ViT} به حافظه‌ی \lr{GPU} بالاتر، مکانیزم \lr{Self-Attention} است. در آن، ماتریس توجه برای $N$ توکن (پچ) دارای پیچیدگی حافظه‌ی $O(N^2)$ است؛ چون باید شباهت هر پچ با همه‌ی پچ‌های دیگر محاسبه و ذخیره شود. در مقابل، کانولوشن در \lr{CNN} ماهیتی محلی با اشتراک وزن دارد و مصرف حافظه‌اش به‌صورت خطی با اندازه‌ی ورودی رشد می‌کند. این رشد درجه‌دوم نقشه‌ی توجه نسبت به تعداد توکن‌ها، عامل تعیین‌کننده‌ی مصرف حافظه‌ی بالاتر \lr{ViT} است.